\documentclass[11pt]{article}
\usepackage[margin=1in]{geometry}
\usepackage[T1]{fontenc}
\usepackage{lmodern,amsmath,amssymb,booktabs,graphicx}
\usepackage[colorlinks=true,linkcolor=blue,urlcolor=blue]{hyperref}
\setlength{\parindent}{0pt}\setlength{\parskip}{5pt}
\title{RoCE-K: continuous-covariate linear-balance validation}
\author{}\date{4 October 2026}
\begin{document}\maketitle

\textbf{Finding.} A feasible conditional linear-model branch now extends the
finite-stratum benchmark to continuous covariates. Exact signed balancing
removes the common linear outcome fit, and leave-out variances remove
sampling noise from source dispersion. It improves interval length against
an equally protected target-only reference, including a $p=200$ check, but
remains conservative relative to ordinary target-normal intervals. This is
not a validated general nonlinear or sparse $p>n$ RoCE procedure.

\section{What was proved and implemented}

Every site's conditional outcome mean, including invalid sources, must lie
in the chosen linear basis. Invalid coefficients may differ arbitrarily
within the bounded-outcome model; no minimum deviation is imposed. At least
$3K/4$ sources are valid in each arm, with possibly different valid sets.
QR computes minimum-norm signed weights matching the target empirical feature
means. All 1,000 observations per site are used; no true nuisance parameters
are supplied to inference.

Within this branch, any common linear fitted predictor cancels algebraically.
The leave-out variance correction is unbiased under heteroskedasticity. The
source dispersion statistic becomes a zero-diagonal patient quadratic form,
for which the main design supplies a finite-sample conditional MGF bound and
an explicit upper confidence bound. It requires no simultaneous confidence
interval for every source variance and no Gaussian approximation to source
estimates. These are new derivations for this benchmark, built on established
balancing and leave-out principles.

Arm-wise and contrast dispersion certificates are compared using a simultaneous
error allocation. The resulting source band is combined with the full-target
conditional band; target-covariate composition uncertainty is added once.
The primary population analysis uses the universal CATE range bound 2.
A separate range-.5 sensitivity assumes that stronger model class in advance.
It must not be presented as a data-estimated range.

A target projection supplies the recorded point estimate; the interval midpoint
is also evaluated. Rank or leverage problems trigger a design-only protected
fallback. There were no design failures in this study. This use of all outcomes
is justified by design-only weights and exact linear cancellation; it does not
justify ignoring cross-fitting boundaries in the general fitted RoCE method.

\section{Prespecified experiments and checks}

There are 6,600 distinct setting-specific repetitions, each with fresh covariates,
treatments and outcomes: 18 main settings, nine boundary settings and six
weak-bias stress settings, all with 200 repetitions. Seven method rows per
setting give 231 summaries. A 15-repetition runtime pilot and 18 same-seed
replays are excluded from the reported Monte Carlo count.

The target has independent uniform $[-1,1]$ covariates. Source covariates use
density $2^{-p}\prod_{r=1}^4(1+\delta_r x_r)$, with
$(\delta_1,\ldots,\delta_4)=s(1,-.6,.4,-.3)$ and default $s=.4$.
Shift sensitivities use $s=0,.8$. Target treatment probabilities are
$.5+.12x_1-.08x_2$; source $j$ adds $.04\cos(j-1)$. The target outcome means are
\[
 m_t^0(x)=.35+.06x_1-.035x_2+.025x_3+.02x_4,\qquad
 m_t^1(x)=m_t^0(x)+.10+h x_1.
\]
Thus the population TATE is fixed at .10. Main and boundary panels use
$h=.15$; the weak-bias stress panel uses $h=0$. The main panel has
$p\in\{4,10,20\}$ and $K\in\{16,64,256\}$ under all-valid and weak-bias
configurations. The boundary panel adds $p=50,100,200$ at $K=64$, source-tilt
sensitivities, strong shifts, cancellation, disjoint arm-valid sets and deliberate
nonlinearity. Weak treated-arm shifts affect one quarter of sources and equal
$2\sqrt{.5/1000}\,K^{-1/4}$; strong shifts equal .15.

All feasible intervals use observed summaries and fixed class assumptions.
Oracle conditional means/variances are used only in audits. Across in-model
settings, maximum valid-source drift was $1.84\times10^{-15}$; maximum balance
error was $1.56\times10^{-15}$. Mean estimated/oracle source variance ratios
ranged from .9990 to 1.0009. No dispersion upper-bound failure was observed.
The full research suite passed 51 tests, including dense-matrix checks, 256
exact non-Gaussian configurations, literal leave-one-out refits and rank failure.
Representative saved intervals reproduced exactly.

\section{Continuous and moderately high-dimensional results}

The following are weak-bias results for the adaptive dispersion procedure with
the universal range bound. All comparators use the same patient samples.
The last column is a paired length difference relative to protected target-only,
with its Monte Carlo standard error.

\begin{center}\small
\resizebox{\linewidth}{!}{\input{continuous_table}}
\end{center}

At $p=20,K=256$, length falls from .3792 to .2623, approximately 31\%.
At $p=200,K=64$, it falls from .4081 to .3485, approximately 15\%.
The $p=200$ regression remains below the roughly 500 observations per treatment
arm; no sparse $p>n$ claim follows. Signed weighting and leverage also become
more costly as dimension increases.

Ordinary target-normal intervals are substantially shorter. Protected coverage
ranges from .995 to 1 over the in-model panels, so conservatism remains clear.
With 200 repetitions, zero observed failures has a two-sided exact 95\%
coverage lower limit of about .982. Observed coverage alone is not the proof.

\begin{figure}[ht]\centering
\includegraphics[width=\linewidth]{continuous_comparison.pdf}
\caption{Weak-bias interval lengths relative to protected target-only (left)
and ordinary target-normal (right), using the universal CATE range bound.
Each setting has 200 fresh repetitions and 1,000 patients per site.}
\end{figure}

\section{Weak bias with growing source count}

The stress panel has $p=10$ and constant target treatment effect. Consequently
its conditional target effect equals its population effect; the naive conditional
source interval and the protected population interval cover the same truth.
The naive interval uses estimated leave-out variances and ignores source bias.

\begin{center}\small
\resizebox{\linewidth}{!}{\input{continuous_stress_table}}
\end{center}

The all-valid naive control remains near nominal coverage. Under weak shifts,
coverage falls to .360 at $K=2048$, even though each individual shift decreases.
The protected method retains coverage, with length .2310 versus .3781 for
protected target-only. This confirms the intended stress mechanism rather than
an artifact of deliberately underestimated source variances.

\begin{figure}[ht]\centering
\includegraphics[width=\linewidth]{continuous_stress.pdf}
\caption{Constant-effect weak-bias stress test. Left: naive pooling and protected
coverage. Right: the universal range-2 procedure and the stronger declared
range-.5 sensitivity. The smaller range is an additional assumption, not an
estimated feature of these data.}
\end{figure}

\section{Boundaries and point-estimate tradeoffs}

\begin{center}\small
\resizebox{\linewidth}{!}{\input{boundary_table}}
\end{center}

Projection improves strong-bias RMSE relative to the midpoint in this panel,
while midpoint is substantially better under cancellation and disjoint weak
shifts. No point rule is selected separately for each scenario.
The nonlinear row deliberately violates the all-site linear-mean requirement.
Its good observed coverage does not extend the theorem: valid transported
sources then have conditional fitted drift as large as .0180 in this panel, and the leave-out variance
correction is no longer justified by the preceding identity.

The remaining bottlenecks are now more specific. The universal range-2 target
composition allowance is conservative, especially for nearly constant treatment
effects. A data-derived range would need its own simultaneous certificate.
For sparse or nonlinear outcome fitting, approximate balance leaves a remainder;
that remainder must enter both source validity and quadratic noise correction.
These issues should be solved before labeling this branch a general RoCE-K method.

\end{document}
